\chapter{Introduction}\label{ch:introduction}
\thispagestyle{pagestyle}

\section{Motivation}
Quantum computing is moving from textbook algorithms toward concrete uses in signal and image processing~\cite{Ruan2021OpportunitiesChallenges,Wang2022ReviewQIP,Farooq2025SystematicReview}. The promise is straightforward. Encode an image in the state of a small qubit register, then run an operation such as filtering, segmentation, or edge detection directly on that state~\cite{Le2011FRQI,Yao2017QHED}. On current hardware the promise meets a wall. Noisy intermediate-scale quantum (NISQ) devices have few qubits and short coherence times, and their gate and readout errors are far from negligible. As a result, circuit depth and two-qubit gate count decide whether a pipeline runs at all, long before any asymptotic speedup could matter~\cite{Preskill2018NISQ,Geng2022HybridNISQEdge,Geng2023ImprovedFRQI}.

This thesis studies one concrete pipeline end to end: the flexible representation of quantum images (FRQI)~\cite{Le2011FRQI} followed by the quantum Hadamard edge detector (QHED)~\cite{Yao2017QHED,QiskitTextbookQuantumEdgeDetection}. FRQI is a natural test case because its state preparation is dominated by multi-controlled gates, which are exactly the operations that inflate depth and CX count on NISQ devices. \Cref{ch:related} defines these terms and reviews the surrounding literature.

\section{Problem and research question}
FRQI stores each grayscale intensity in a rotation angle on one \emph{color} qubit, addressed by a \emph{position} register of $m = 2\log_2 N$ qubits for an $N\times N$ image~\cite{Le2011FRQI,Iliyasu2013FRQIroadmap}. Preparing the state means looping over pixel addresses and applying a multi-controlled rotation for each one. The cost of that loop depends entirely on how the multi-controlled gates are synthesized, and modern synthesis work shows that ancilla-assisted decompositions can cut CX depth substantially~\cite{Vale2023MCSU2,Azevedo2024LowDepthMCG,Arsoski2024MCXQFT,Zindorf2024EfficientMCG}.

The research question is deliberately narrow, and it does not propose a new edge operator. What it asks is this: \textbf{when the only change is the multi-controlled synthesis used to prepare the same FRQI state, how do transpiled resources and noise robustness change, and does any resource saving carry through to the edge-detection result?} To keep the comparison honest, QHED and its classical Sobel reference are held fixed; only the preparation circuit and its noise channel vary. A small readout-mitigation slice is added on the color qubit because readout error biases the inferred intensities and is a standard concern on real devices~\cite{Mu2025GuidedImageFiltering}.

Prior quantum edge-detection work usually assumes ideal preparation or studies NEQR-based Sobel variants~\cite{Zhang2013NEQR,Zhang2015QSobel,Fan2019SobelNEQR,Chetia2021ImprovedSobelNEQR}. Few studies compare naive against ancilla-assisted FRQI preparation on identical noise grids while keeping the downstream stage fixed. \Cref{ch:related} positions the present work against that gap.

\section{Contributions}
This thesis makes three empirical contributions and, most importantly, one methodological contribution. All of them are backed by reproducible scripts and CSV manifests in the project repository.\footnote{\url{https://github.com/AlexandruAlex10/qip-frqi-qhed-mcx-ancilla-and-calibration-matrix-inversion}}
\begin{enumerate}
  \item A \textbf{structural comparison} of naive (\texttt{noancilla}) versus v-chain (ancilla-assisted) multi-controlled FRQI preparation, reporting transpiled qubit count, depth, and CX for $4\times 4$, $8\times 8$, and $16\times 16$ images, with explicit scaling estimates where full naive transpilation is infeasible.
  \item A \textbf{noisy-simulator robustness study} that reconstructs FRQI states under depolarizing noise, runs QHED, and reports PSNR, SSIM, fidelity, and edge metrics for both layouts on matched noise grids.
  \item A \textbf{readout-mitigation proof of concept} that estimates a $2\times 2$ confusion matrix for the color qubit, inverts it with regularization, and compares raw versus mitigated $p_1$ across random seeds on a $4\times 4$ slice.
  \item \textbf{(Main contribution) A comparative methodology for NISQ image pipelines:} hold the downstream operator fixed, vary only the preparation synthesis, and report resources, fidelity proxies, and edge-centric metrics \emph{together}. Keeping them in one view prevents a single gate-count headline from hiding an inconclusive application result, which turned out to be the central lesson of the study.
\end{enumerate}

\section{Scope and limitations}
All experiments use \textbf{simulators only} (Qiskit Aer statevector or density-matrix backends, with optional mock NISQ noise). No IBM Quantum hardware runs are included. Images are limited to $4\times 4$, $8\times 8$, and $16\times 16$ grayscale test patterns. Density-matrix simulation of the full $16\times 16$ v-chain register (16~qubits) is often skipped to respect memory limits (insufficient RAM); \Cref{ch:results} marks where v-chain data is missing. The statistical tests between naive and v-chain are reported as \textbf{exploratory} and must not be read as definitive superiority claims.

\section{Thesis outline}
\Cref{ch:related} covers the background this study leans on, from quantum-computing essentials and image metrics through quantum image representations, the edge-detection literature, NISQ constraints, multi-controlled synthesis, and readout mitigation, and it closes by positioning the work. The pipeline itself, together with the preparation modes, noise model, metrics, statistics, and experimental setup, is formalized in \Cref{ch:methods}. \Cref{ch:results} then walks through the figures and tables and interprets them, and \Cref{ch:conclusions} draws the contributions and their implications together before pointing to future work. Supporting material, including the symbol tables, a reproduction checklist, the repository and architecture notes, and the extended result tables, sits in the appendices.
